\documentclass[11pt]{article}
\usepackage[margin=1in]{geometry}
\usepackage[T1]{fontenc}
\usepackage{lmodern}
\usepackage{textcomp}
\usepackage{amsmath,amssymb,amsthm}
\usepackage{booktabs}
\usepackage{pgfplots}
\pgfplotsset{compat=1.17}
\usepackage[hidelinks]{hyperref}
\usepackage{enumitem}
\setlist{nosep,leftmargin=*}
\definecolor{seriesA}{HTML}{2A78D6}
\definecolor{seriesB}{HTML}{EB6834}
\definecolor{inkMuted}{HTML}{8A8A85}
\newtheorem{theorem}{Theorem}
\newtheorem{definition}{Definition}
\newtheorem{conjecture}{Conjecture}
\theoremstyle{remark}
\newtheorem*{remark}{Remark}
\newtheorem{question}{Question}
\newcommand{\lean}[1]{\texttt{#1}}
\emergencystretch=3em

\title{Fruitfulness as Compression:\\
Lean-Checked Theorems and a Measurement of All of Mathlib}
\author{Vasily Shablov\thanks{Formalization, measurement and drafting assisted by Claude (Anthropic).}}
\date{October 2026 \\ \small Draft}

\begin{document}
\maketitle

\begin{abstract}
We pose an open-ended question: what is the \emph{fruitfulness} of a mathematical idea?
We answer it with a definition: a named concept is a compression, and its value is what
naming it saves. We then prove theorems about this definition in Lean~4 and measure it on
all 771{,}129 constants of Mathlib.

\textbf{Theorems.}
\begin{itemize}
\item The saving from naming a body of size $b$ used $k$ times is exactly $(k-1)(b-1)-1$.
\item Unfolding all names at most doubles proof size per result, and the doubling is
  attained.
\item With at most $d$ citations per result, the sharp bound is the $d$-bonacci recurrence,
  also attained.
\end{itemize}

\textbf{Measurements.}
\begin{enumerate}
\item \emph{Growth law.} Naively unfolded proof size grows by a factor of about $1.5$ per
  dependency level ($R^2=0.98$), reaching about $10^{65.6}$ symbols.
\item \emph{Rediscovery.} Ranking by reuse alone rediscovers the textbook core:
  \lean{Set}, \lean{Category}, \lean{Functor}, \lean{TopologicalSpace}, \lean{Real},
  \lean{Module}, \ldots
\item \emph{Two axes of taste.} Human curation of \emph{concepts} is predicted by reuse
  (AUC $0.88$). Human fame of \emph{theorems} is predicted by depth (AUC $0.78$) and not
  by reuse (AUC $0.54$).
\item \emph{A glue layer.} Instance declarations and structure projections are $12\%$ of
  constants but receive $60\%$ of all citations. Mathlib also obeys Zipf's law of
  abbreviation. These are the same signatures as function words in language and currency
  metabolites in metabolism.
\end{enumerate}

We correct an overreach in our own earlier statement (naive unfolding is not minimal
lemma-free proof size). We connect the work to proof complexity, logical depth,
compression progress, hierarchy in nature, and Wigner's puzzle. We state nine open
questions and five falsifiable conjectures.
\end{abstract}

\section{Introduction}

\paragraph{Where the question came from.} The question started as a debate. One position
holds that AI is good at well-posed yes/no problems but cannot handle open ``why / what
is'' questions, because ``there is nothing to optimize.'' Rather than argue, we posed one
such question and followed it:

\begin{quote}\itshape
What is the fruitfulness of a mathematical idea, and can it be measured?
\end{quote}

\paragraph{Our answer.} We take a compression view, in the tradition of minimum
description length and of Whitehead's remark that ``civilization advances by extending the
number of important operations which we can perform without thinking about them.''
\begin{itemize}
\item A definition or lemma is \emph{named}: its body is written once and then cited.
\item The value of a name is the description length it saves.
\item This makes fruitfulness precise enough to prove theorems about, and concrete enough
  to measure on the largest library of formalized mathematics.
\end{itemize}

\paragraph{Contributions.}
\begin{enumerate}
\item \textbf{Theorems, checked in Lean 4 with Mathlib} (Section~\ref{sec:theory}; no
  \lean{sorry}; standard axioms only).
  \begin{itemize}
  \item The exact value of naming a concept.
  \item A tight exponential bound on unfolding.
  \item A sharp bound for bounded citation, via the $d$-bonacci recurrence.
  \end{itemize}
\item \textbf{A complete measurement of Mathlib} (Sections~\ref{sec:data}--\ref{sec:results}).
  \begin{itemize}
  \item A growth law for unfolded size.
  \item Reuse rediscovers the textbook core.
  \item Two distinct axes of human mathematical taste, each predicted by a computable
    measure.
  \item A glue layer with the statistics of function words.
  \end{itemize}
\item \textbf{Connections and open questions} (Section~\ref{sec:questions}).
  \begin{itemize}
  \item The deepest version of ``are concepts necessary?'' is the open Frege versus
    Extended Frege problem.
  \item Human fame tracks Bennett's logical depth; concept taste tracks Schmidhuber's
    compression progress.
  \item The same reuse statistics appear in language, metabolism and software.
  \end{itemize}
\end{enumerate}

\section{Theory: what naming buys}\label{sec:theory}

\begin{definition}[Library]
A \emph{library} consists of:
\begin{itemize}
\item a size function $\mathrm{size}:\mathbb{N}\to\mathbb{N}$;
\item citation sets $\mathrm{deps}(i)\subseteq\{0,\dots,i-1\}$.
\end{itemize}
Result $i$ has its own proof of size $\mathrm{size}(i)$ and cites the results in
$\mathrm{deps}(i)$. Its \emph{unfolded size} is
\[
U(i)=\mathrm{size}(i)+\sum_{j\in\mathrm{deps}(i)}U(j),
\]
which is the size of its proof if each cited result were re-proved in place.
\end{definition}

\begin{theorem}[Value of one name; \lean{naming\_saves}, \lean{naming\_pays\_iff}]
Inlining a body of size $b$ at $k$ sites costs $kb$. Naming it costs $b+k$: one
definition plus one reference per use. The saving is exactly
\[
kb-(b+k)=(k-1)(b-1)-1.
\]
\begin{itemize}
\item Naming pays off if and only if $(k-1)(b-1)>1$. In particular it always pays when
  $k\ge 2$ and $b\ge 3$ (\lean{naming\_pays\_of}).
\item A name used once never pays: it costs exactly one symbol (\lean{naming\_once}).
\end{itemize}
\end{theorem}

\noindent The last clause is the \emph{rule utility} constraint of the Sequitur grammar
compressor~\cite{sequitur}, here derived rather than imposed.

\begin{theorem}[Doubling; \lean{Library.unfold\_le\_two\_pow}, \lean{completeLib\_unfold}]
If $\mathrm{size}\le S$, then $U(i)\le S\cdot 2^{i}$, whatever the number of citations.
Equality holds for the \emph{complete library}, in which every result cites all earlier
ones and has size $1$: there $U(i)=2^{i}$.
\end{theorem}

\begin{theorem}[Sharp bound for bounded citation; \lean{Library.unfold\_le\_tb}, \lean{windowLib\_unfold}]
Let $T_d(i)=1+\sum_{j=\max(0,i-d)}^{i-1}T_d(j)$. If $\mathrm{size}\le S$ and
$|\mathrm{deps}(i)|\le d$ for all $i$, then
\[
U(i)\le S\cdot T_d(i).
\]
Equality holds for the library in which each result cites its $d$ predecessors and has
size $1$.
\end{theorem}

\noindent The proof rests on an exchange lemma (\lean{sum\_le\_sum\_window}): among sets of
at most $d$ indices below $i$, a monotone function has its largest sum on the last $d$
indices. Some consequences:
\begin{itemize}
\item $T_d(i)$ grows like $\phi_d^{\,i}$, where $\phi_d$ is the $d$-bonacci constant, the
  root in $(1,2)$ of $x^d=x^{d-1}+\dots+1$.
\item $\phi_2$ is the golden ratio. The case $d=2$ is formalized separately: $U(i)=F_{i+3}-1\ge 2^{\lfloor i/2\rfloor}$
  (\lean{fibLib\_unfold}, \lean{fibLib\_unfold\_ge}).
\item $\phi_d\to 2$ as $d\to\infty$, recovering the doubling theorem.
\end{itemize}
Every one of these exponential gaps is generated by a library of \emph{linear} named size.

\begin{remark}[What these theorems do and do not say]
$U$ measures \emph{copying}. It is not the size of the shortest lemma-free proof, and a
clever re-prover can restructure a proof rather than inline it. For Frege systems,
tree-like proofs polynomially simulate dag-like ones~\cite{krajicek}. So, in that setting,
reusing lemmas saves only polynomially, even though copying costs exponentially. The
question of whether \emph{new symbols} (definitions) can save super-polynomially is the
Frege versus Extended Frege problem, which is open (Section~\ref{sec:questions}). Our
theorems are about naming versus copying, and they are tight for that.
\end{remark}

\section{Data and method}\label{sec:data}

\paragraph{Export.} We loaded all of Mathlib~\cite{mathlib} (Lean v4.33.1), including private proof
bodies, and recorded every constant: its kind, its module, the size of its body (number of
distinct term nodes), and the constants cited in its body and statement. In a second pass,
with environment extensions loaded, we recorded the size of its statement and whether it is
an instance declaration or a structure projection.

\paragraph{Size of the data.}
\begin{itemize}
\item $771{,}129$ constants and $20.4$\,M citation edges.
\item $279$ cycle edges, from unsafe definitions, were dropped.
\item The \emph{visible} set has $315{,}085$ constants: user-facing Mathlib theorems,
  definitions and inductive types, after removing compiler-generated auxiliaries.
\end{itemize}

\paragraph{Measures.} For each constant $c$:
\begin{itemize}
\item \emph{reuse} $k(c)$: the number of constants citing $c$;
\item \emph{unfolded size} $U(c)$: as in Section~\ref{sec:theory}, computed exactly in
  log-space;
\item \emph{height}: the longest citation chain below $c$;
\item \emph{depth}: $\log_{10}U(c)-\log_{10}(\text{statement size})$, i.e.\ hidden work
  per symbol of statement.
\end{itemize}

\paragraph{Glue.} We call instance declarations and structure projections \emph{glue},
classified structurally rather than by name. They make up $12.3\%$ of visible constants.

\paragraph{Human labels.} All labels come from Mathlib's own human-curated files.
\begin{itemize}
\item \emph{Famous}: Freek Wiedijk's ``100 theorems'' and the ``1000+ theorems'' project,
  294 declarations.
\item \emph{Curriculum}: the undergraduate curriculum list and the library overview, 686
  declarations.
\end{itemize}

\paragraph{Statistics.}
\begin{itemize}
\item AUC against all other non-glue constants of the same kind. $0.5$ is chance and
  $1.0$ is perfect.
\item A \emph{same-file} control: the probability that a labelled constant beats an
  unlabelled constant of the same kind in the same source file. This controls for
  ``advanced files contain deeper things.''
\end{itemize}

\section{Results}\label{sec:results}

\subsection{A growth law}

Mean $\log_{10}U$ is almost exactly linear in height (Figure~\ref{fig:growth}).
\begin{itemize}
\item \textbf{Fit.} Over all $315{,}085$ constants, $\log_{10}U\approx1.22+0.189\,h$
  with $R^2=0.979$, a factor of $1.55$ per level. Fitting the per-height means gives
  $1.50$ per level with $R^2=0.990$.
\item \textbf{Stability across the range.} The three thirds of the height range give
  $1.45$, $1.64$ and $1.41$.
\item \textbf{Placement.} This is supercritical branching: well above the linear regime,
  and below the golden ratio, $1.618$ per \emph{result}, of the citing-two-predecessors
  library.
\item \textbf{Scale.} The deepest result, \lean{CompletelyPositiveMap.map\_cstarMatrix\_nonneg}
  at height 361, has $U\approx10^{65.6}$.
\end{itemize}

\begin{figure}[t]
\centering
\begin{tikzpicture}
\begin{axis}[width=0.92\linewidth,height=6.2cm,
  xlabel={dependency height $h$}, ylabel={$\log_{10}$ unfolded size $U$},
  xmin=0,xmax=365,ymin=0,ymax=70,
  axis lines=left, grid=major, grid style={draw=black!8},
  tick label style={font=\small,text=black!70}, label style={font=\small},
  legend style={draw=none,font=\small,at={(0.03,0.97)},anchor=north west,fill=none},
  legend cell align=left]
\addplot[color=seriesB,line width=1pt] table[x=h,y=max]{../data/fig_growth.dat};
\addlegendentry{maximum at each height}
\addplot[color=seriesA,line width=1.2pt] table[x=h,y=mean]{../data/fig_growth.dat};
\addlegendentry{mean at each height}
\addplot[color=inkMuted,dashed,line width=0.8pt,domain=0:362]{1.216+0.18898*x};
\addlegendentry{fit $1.22+0.189h$ ($R^2=0.98$)}
\end{axis}
\end{tikzpicture}
\caption{Naively unfolded size grows exponentially with dependency height, about
$1.5\times$ per level, over $315{,}085$ Mathlib constants.}
\label{fig:growth}
\end{figure}

\subsection{Reuse rediscovers the textbook}

With glue removed, the most-cited definitions and theorems are listed below. No notion of
importance was supplied. The ranking is pure citation counting.

\begin{center}\small
\begin{tabular}{p{0.46\linewidth}p{0.46\linewidth}}
\toprule
\textbf{Definitions} (by citations) & \textbf{Theorems} (by citations) \\
\midrule
\lean{Set}, \lean{Category}, \lean{Functor}, \lean{TopologicalSpace}, \lean{Real},
\lean{Module}, \lean{RingHom.id}, \lean{CommRing}, \lean{Finset}, \lean{Semiring},
\lean{RingHom}, \lean{AddCommMonoid}, \lean{AddCommGroup}, \lean{Algebra},
\lean{NormedAddCommGroup}, \lean{LinearMap}, \lean{Equiv}, \lean{Opposite}, \ldots
&
\lean{mul\_one}, \lean{add\_zero}, \lean{le\_trans}, \lean{zero\_add}, \lean{one\_mul},
\lean{Nat.cast\_one}, \lean{mul\_comm}, \lean{Finset.sum\_congr}, \lean{Set.ext},
\lean{le\_refl}, \lean{mul\_assoc}, \lean{map\_zero}, \lean{le\_antisymm}, \ldots \\
\bottomrule
\end{tabular}
\end{center}

\noindent The theorem column omits tactic-internal certificate lemmas
(\lean{Mathlib.Tactic.*}), which form a third, automation-generated kind of glue.

\paragraph{Observation.} The top of the concept ranking is dominated by notions of
\emph{structure and sameness}: sets, maps, homomorphisms, equivalences, categories,
functors, opposites. The single most-cited definition overall is \lean{DFunLike.coe}
(applying a structured map), at $115{,}373$ citations. In a precise sense, the most
compressive idea in modern mathematics is the morphism.

\subsection{Two axes of human taste}

Table~\ref{tab:auc} gives the main result.
\begin{itemize}
\item \textbf{Concepts.} Human-curated concepts are predicted by \emph{reuse}: AUC $0.88$,
  and $0.87$ against peers in the same file. Depth does not predict them, and is even
  inverted (AUC $0.46$): fundamental concepts are shallow.
\item \textbf{Theorems.} Famous and curriculum theorems are predicted by \emph{depth}:
  AUC $0.78$--$0.82$, and $0.66$--$0.68$ within the same file. Reuse barely predicts them
  ($0.54$--$0.58$).
\item \textbf{Famous theorems are not reuse hubs.} $28.6\%$ of famous declarations are never
  cited anywhere else in Mathlib, against $36.6\%$ of all theorems, so they are only
  slightly more reused than average.
\end{itemize}

\begin{table}[t]
\centering\small
\begin{tabular}{lrcccc}
\toprule
Human label & $n$ & reuse $k$ & unfolded $U$ & depth & statement size \\
\midrule
curriculum \emph{definitions} & 421 & \textbf{0.88} / \textbf{0.87} & 0.46 / 0.35 & 0.46 / 0.38 & 0.36 / 0.32 \\
curriculum theorems & 209 & 0.58 / 0.59 & \textbf{0.82} / 0.66 & \textbf{0.82} / \textbf{0.68} & 0.67 / 0.51 \\
famous theorems & 259 & 0.54 / 0.55 & \textbf{0.78} / 0.66 & \textbf{0.78} / \textbf{0.67} & 0.59 / 0.51 \\
\bottomrule
\end{tabular}
\caption{Global AUC / same-file win rate, glue removed. $0.5$ is chance.}
\label{tab:auc}
\end{table}

\subsection{A glue layer, a heavy tail, and the law of abbreviation}

\paragraph{Glue.} Glue (instances and projections) is $12.3\%$ of visible constants but
receives $59.5\%$ of all citations, including 77 of the 100 most-cited constants.
\begin{itemize}
\item Only 63 of the 904 human-labelled declarations are glue.
\item This is the profile of \emph{function words} in language: a small closed class
  carrying a large share of tokens and little content.
\item It is also the profile of \emph{currency metabolites} (ATP, NADH, water), which
  network biologists must remove before metabolic pathways become visible~\cite{ma}.
\end{itemize}

\paragraph{Heavy tail.} Reuse is heavy-tailed (Figure~\ref{fig:ccdf}).
\begin{itemize}
\item A discrete power-law maximum-likelihood fit~\cite{clauset} gives a density exponent
  $\alpha\approx1.72$, $1.77$ and $1.79$ for $k\ge10$, $30$ and $100$.
\item For comparison: word frequencies give $\alpha\approx2$ (Zipf), and protein-domain
  family sizes and citation counts show similar tails~\cite{koonin,price}.
\end{itemize}

\paragraph{Law of abbreviation.} In natural languages, frequent words are shorter (Zipf),
a regularity proposed to be a universal consequence of compression~\cite{ferrer}. Mathlib
obeys it too.
\begin{itemize}
\item \textbf{Definitions.} Mean name length falls monotonically with reuse: $13.2$
  characters (never cited), $10.8$ (cited $10$--$99$ times), $8.9$ ($10^3$--$10^4$),
  $8.2$ ($10^4$--$10^5$), and $3$ for \lean{Set}.
\item \textbf{Theorems.} Mean name length falls from $19.2$ to $8.5$.
\item \textbf{Rank correlations.} Spearman $\rho(k,\text{name length})=-0.16$ for
  definitions and $-0.07$ for theorems; $\rho(k,\text{body size})=-0.22$ for definitions.
  Most constants have small $k$, so the global correlations are weak, but the trend
  across decades of $k$ is monotone.
\end{itemize}
Mathlib's names are derived from statements by convention, so this is not designed; it
emerges from which concepts get reused.

\begin{figure}[t]
\centering
\begin{tikzpicture}
\begin{loglogaxis}[width=0.72\linewidth,height=5.6cm,
  xlabel={number of citations $x$}, ylabel={constants with $k\ge x$},
  axis lines=left, grid=major, grid style={draw=black!8},
  tick label style={font=\small,text=black!70}, label style={font=\small}]
\addplot[color=seriesA,line width=1.2pt,mark=*,mark size=1.2pt] table[x=k,y=ccdf]{../data/fig_ccdf.dat};
\end{loglogaxis}
\end{tikzpicture}
\caption{Reuse in Mathlib is heavy-tailed: 215{,}078 constants are cited at least once,
8{,}090 at least 100 times, and 165 at least 10{,}000 times.}
\label{fig:ccdf}
\end{figure}

\section{Connections and open questions}\label{sec:questions}

\subsection{Are concepts necessary? Proof complexity}

Naming comes in two strengths.
\begin{itemize}
\item \textbf{Lemmas} reuse proved statements: dag-like versus tree-like proofs.
  \begin{itemize}
  \item The gap is exponential for resolution~\cite{bsiw}.
  \item It is polynomial for Frege~\cite{krajicek}.
  \end{itemize}
\item \textbf{Definitions} introduce new symbols. This is Tseitin's extension rule; Frege
  plus extension is Extended Frege~\cite{cookreckhow}.
  \begin{itemize}
  \item Whether Extended Frege is super-polynomially stronger than Frege is a central open
    problem.
  \end{itemize}
\end{itemize}
Higher up, the answer is known to be dramatic. G\"odel~\cite{godel} showed that passing to
higher-order logic shortens some proofs by more than any computable function, and Boolos's
``curious inference''~\cite{boolos} gives a concrete first-order-valid argument whose
first-order proofs are physically unwritable but which has a short second-order proof.

\begin{question}[Are new concepts forced?]
Are \emph{definitions} ever super-polynomially necessary beyond lemma reuse? This is
Frege versus Extended Frege, restated as the question ``are human concepts forced or merely
convenient?''
\end{question}

\begin{question}[Empirical proof complexity]
Mathlib is a vast sample of definitions in actual use. Can it locate \emph{where}
definitions give large savings, and so supply empirical evidence on a problem that has
resisted proof?
\end{question}

\subsection{Value has two axes: depth and leverage}

The split in Table~\ref{tab:auc} matches two independent theories of value.
\begin{itemize}
\item \textbf{Bennett's logical depth}~\cite{bennett}. An object is valuable when it is
  the product of a long computation that cannot be shortcut. Famous theorems are deep: a
  short statement certifies a large amount of hidden work.
\item \textbf{Schmidhuber's compression progress}~\cite{schmidhuber}. Interest equals
  improvement of the observer's compressor. Our $(k-1)(b-1)-1$ is exactly the compression
  progress of one new concept, and reuse predicts concept taste.
\end{itemize}
We know of no previous large-scale test of either thesis against human judgments of
mathematical value, though a fuller literature search is warranted.

\begin{question}[The residual]
Are depth and leverage the only two axes of mathematical value? What explains the
remaining $12$--$22\%$ of AUC: surprise, beauty, \emph{bridging} between distant areas, or
historical accident?
\end{question}

\begin{question}[Computability of taste]
Ideal fruitfulness, meaning the optimal library of concepts for all of mathematics, is a
minimal-grammar problem. It is NP-hard even for finite corpora~\cite{charikar}, and in the
limit it is Kolmogorov complexity, which is uncomputable. So no agent, human or machine,
can optimize fruitfulness exactly; all can only approximate it from below. Is the human
advantage in choosing directions, if any, then a matter of better approximation rather than
a different kind of faculty?
\end{question}

\subsection{A glue law across nature}

Function words, currency metabolites, tiny utility packages (the \lean{left-pad} incident),
and Mathlib's instances share a profile: few units, an enormous share of uses, little
content, and invisible to anyone asked what matters.

\begin{question}[The glue theorem]
Must every system that grows by reuse under a description-length cost develop (i)
heavy-tailed reuse, (ii) a high-frequency, low-content glue layer, and (iii) abbreviation of
the frequent? Preferential attachment~\cite{price} explains (i), and compression
arguments~\cite{ferrer} explain (iii). Is there a single derivation of all three?
\end{question}

\noindent Mathlib is a natural \emph{model organism} for such laws.
\begin{itemize}
\item In biology, edges are inferred and noisy. In language, ``concept'' is fuzzy.
\item In Mathlib, every edge is exact and every unit is machine-checked.
\end{itemize}

\begin{question}[Mathematics as an evolving genome]
Protein-domain statistics are explained by birth--death--innovation models~\cite{koonin}.
Does the same model, fitted to Mathlib's history (generalization as duplication, imported
concepts as horizontal transfer), reproduce its reuse distribution?
\end{question}

\subsection{Hierarchy, decoupling, and the effectiveness of mathematics}

\paragraph{Simon's watchmakers.} Simon's parable~\cite{simon} contrasts a watchmaker who
builds from stable sub-assemblies with one who does not. The first has an exponential
advantage under interruption, which Simon offered as the reason nature is hierarchical.
Our theorems are the corresponding statement about \emph{size}.

\paragraph{Proof irrelevance is scale separation.} In Lean, any two proofs of the same
proposition are equal. Later results therefore depend only on \emph{statements}, never on
proofs. This is the logical form of the physicist's decoupling of scales: chemistry needs
the periodic table, not quantum chromodynamics. The depth $\log U-\log(\text{statement})$
measures how much micro-detail a statement forgets.

\begin{question}[Wigner's puzzle as compression]
Is the ``unreasonable effectiveness'' of mathematics in physics~\cite{wigner} a consequence
of both being hierarchical compressions under scale separation? Formalized physics
libraries now exist. Do they reuse the same top concepts as Mathlib, with the same tail
exponent and glue layer?
\end{question}

\subsection{Delayed reward and understanding}

\begin{question}[How delayed is the reward?]
The strongest form of the ``nothing to optimize'' objection is that fruitfulness is
visible only decades later; Lobachevsky is the canonical example. Mathlib's version history
records when each declaration and each citation appeared. Does early reuse, after months,
predict late reuse, after years? If it does, the reward horizon is short and learnable. If
it does not, the objection gains real empirical support.
\end{question}

\begin{question}[Understanding as conditional compression]
Thurston held that mathematics is about human understanding~\cite{thurston}. Define: an
agent \emph{understands} a proof when its size, relative to the agent's library of named
concepts, fits within the agent's working capacity. This is conditional description length
with a library in place of a universal machine, consistent with chunking in
cognition~\cite{chase,miller}. On this definition, ``can humans understand machine
mathematics?'' becomes the quantitative question ``can the machine's library be taught?''
\end{question}

\subsection{Falsifiable conjectures}

\begin{conjecture}[Growth law]
In future versions of Mathlib, the per-level unfolding factor stays in $[1.4,1.7]$ with
$R^2>0.95$.
\end{conjecture}

\begin{conjecture}[Two axes]
In other large formal libraries (Isabelle AFP, Rocq/MathComp), human-curated concepts are
predicted by reuse and famous theorems by depth, each with AUC $>0.75$. Reuse does not
predict fame (AUC $<0.6$).
\end{conjecture}

\begin{conjecture}[Glue share]
In any mature formal library with typeclass-style inference, a structurally defined glue
class of at most $15\%$ of constants receives a majority of citations.
\end{conjecture}

\begin{conjecture}[Concepts are rediscovered]
Strip names from Mathlib's proof terms and run compression-driven library
learning~\cite{dreamcoder,stitch}. The highest-value recovered abstractions substantially
overlap the top concepts of Section~\ref{sec:results}.
\end{conjecture}

\begin{conjecture}[Morphisms are the universal compressor]
In every sufficiently large formal library, the most-cited definitions are dominated by
notions of map and sameness (morphism, homomorphism, equivalence, functor) rather than by
specific objects.
\end{conjecture}

\section{Limitations}

\begin{itemize}
\item \textbf{The labels are biased.} They record what formalizers chose to curate.
\item \textbf{$U$ is a copying measure.} It counts each citation once per citing result.
  It is not an occurrence count inside proof terms, and not a lower bound on proof size.
\item \textbf{The glue classification is incomplete.} It is structural but leaves out
  tactic-certificate lemmas, which we report separately.
\item \textbf{One snapshot.} All measurements come from a single version of Mathlib.
  Temporal claims are stated as conjectures.
\item \textbf{The reported correlations are observational.}
\end{itemize}

\section{Conclusion}

An open ``what is'' question led us to:
\begin{itemize}
\item a definition;
\item theorems proved in Lean, including a sharp one;
\item a growth law over 315{,}085 objects;
\item a mechanical measure that recovers human concept taste at AUC $0.88$ and theorem
  taste at AUC $0.78$;
\item connections to proof complexity, information theory, biology, linguistics and
  physics;
\item questions we cannot yet answer.
\end{itemize}
Fruitfulness is not unmeasurable. It is delayed and, in its ideal form, uncomputable, for
every kind of mind.

\appendix
\section{Lean index}
All results are in \lean{Fruitfulness.lean}, namespace \lean{Fruitfulness}. They use only
\lean{propext}, \lean{Classical.choice} and \lean{Quot.sound}, with no \lean{sorry}.

\begin{center}\small
\begin{tabular}{ll}
\toprule
Statement & Lean name \\
\midrule
Saving $(k-1)(b-1)-1$ & \lean{naming\_saves} \\
Pays iff $(k-1)(b-1)>1$; pays when $k\ge2$, $b\ge3$ & \lean{naming\_pays\_iff}, \lean{naming\_pays\_of} \\
Used once never pays & \lean{naming\_once} \\
Unfolded size $\ge$ own size & \lean{Library.size\_le\_unfold} \\
$U(i)\le S(d+1)^{i+1}$ & \lean{Library.unfold\_le} \\
$U(i)\le S\,2^i$ (any $d$) & \lean{Library.unfold\_le\_two\_pow} \\
Complete library: $U(i)=2^i$ & \lean{completeLib\_unfold} \\
Exchange lemma & \lean{sum\_le\_sum\_window} \\
$T_d$ is monotone & \lean{tb\_mono} \\
Sharp bound $U(i)\le S\,T_d(i)$ & \lean{Library.unfold\_le\_tb} \\
Attained by the window library & \lean{windowLib\_unfold} \\
Fibonacci library: $U(i)=F_{i+3}-1\ge2^{\lfloor i/2\rfloor}$ & \lean{fibLib\_unfold}, \lean{fibLib\_unfold\_ge} \\
\bottomrule
\end{tabular}
\end{center}

\noindent Reproduction: \lean{Export.lean} and \lean{Flags.lean} export the graph;
\lean{analyze.py}, \lean{robustness.py} and \lean{paper\_stats.py} compute every number in
this paper.

\begin{thebibliography}{99}\small
\bibitem{bsiw} E.~Ben-Sasson, R.~Impagliazzo, A.~Wigderson. Near optimal separation of tree-like and general resolution. \emph{Combinatorica} 24 (2004).
\bibitem{bennett} C.~H.~Bennett. Logical depth and physical complexity. In \emph{The Universal Turing Machine: A Half-Century Survey}, 1988.
\bibitem{boolos} G.~Boolos. A curious inference. \emph{J.\ Philosophical Logic} 16 (1987).
\bibitem{stitch} M.~Bowers et al. Top-down synthesis for library learning. \emph{POPL} 2023.
\bibitem{charikar} M.~Charikar et al. The smallest grammar problem. \emph{IEEE Trans.\ Inf.\ Theory} 51 (2005).
\bibitem{chase} W.~G.~Chase, H.~A.~Simon. Perception in chess. \emph{Cognitive Psychology} 4 (1973).
\bibitem{clauset} A.~Clauset, C.~R.~Shalizi, M.~E.~J.~Newman. Power-law distributions in empirical data. \emph{SIAM Review} 51 (2009).
\bibitem{cookreckhow} S.~A.~Cook, R.~A.~Reckhow. The relative efficiency of propositional proof systems. \emph{J.\ Symbolic Logic} 44 (1979).
\bibitem{dreamcoder} K.~Ellis et al. DreamCoder: bootstrapping inductive program synthesis with wake-sleep library learning. \emph{PLDI} 2021.
\bibitem{ferrer} R.~Ferrer-i-Cancho et al. Compression as a universal principle of animal behavior. \emph{Cognitive Science} 37 (2013).
\bibitem{godel} K.~G\"odel. \"Uber die L\"ange von Beweisen. \emph{Ergebnisse eines math.\ Kolloquiums} 7 (1936).
\bibitem{koonin} E.~V.~Koonin, Y.~I.~Wolf, G.~P.~Karev. The structure of the protein universe and genome evolution. \emph{Nature} 420 (2002).
\bibitem{krajicek} J.~Kraj\'i\v{c}ek. \emph{Bounded Arithmetic, Propositional Logic, and Complexity Theory}. Cambridge University Press, 1995.
\bibitem{ma} H.~Ma, A.-P.~Zeng. Reconstruction of metabolic networks from genome data and analysis of their global structure. \emph{Bioinformatics} 19 (2003).
\bibitem{mathlib} The mathlib Community. The Lean mathematical library. \emph{CPP} 2020.
\bibitem{miller} G.~A.~Miller. The magical number seven, plus or minus two. \emph{Psychological Review} 63 (1956).
\bibitem{sequitur} C.~G.~Nevill-Manning, I.~H.~Witten. Identifying hierarchical structure in sequences: a linear-time algorithm. \emph{JAIR} 7 (1997).
\bibitem{price} D.~de~Solla Price. A general theory of bibliometric and other cumulative advantage processes. \emph{JASIS} 27 (1976).
\bibitem{schmidhuber} J.~Schmidhuber. Driven by compression progress. In \emph{Anticipatory Behavior in Adaptive Learning Systems}, Springer, 2009.
\bibitem{simon} H.~A.~Simon. The architecture of complexity. \emph{Proc.\ Amer.\ Phil.\ Soc.} 106 (1962).
\bibitem{thurston} W.~P.~Thurston. On proof and progress in mathematics. \emph{Bull.\ AMS} 30 (1994).
\bibitem{wigner} E.~Wigner. The unreasonable effectiveness of mathematics in the natural sciences. \emph{Comm.\ Pure Appl.\ Math.} 13 (1960).
\end{thebibliography}
\end{document}
